% =====================================================================
\begin{thebibliography}{99}
\addcontentsline{toc}{chapter}{Bibliografía}
\markboth{Bibliografía}{}
\small

\subsection*{Exámenes universitarios citados}

\bibitem{tripos13p1}
University of Cambridge, Faculty of Mathematics.
\newblock \emph{Mathematical Tripos Part IA}, Paper~1, 30 de mayo de 2013.
\newblock Pregunta~1C (Vectors and Matrices): teorema de De Moivre y raíces $n$-ésimas.
\newblock \url{https://www.maths.cam.ac.uk/undergrad/pastpapers/files/2013/PaperIA_1.pdf}

\bibitem{tripos19p1}
University of Cambridge, Faculty of Mathematics.
\newblock \emph{Mathematical Tripos Part IA}, Paper~1, 30 de mayo de 2019.
\newblock Pregunta~1C(b) (Vectors and Matrices): cociente de potencias de $1\pm i$.
\newblock \url{https://www.maths.cam.ac.uk/undergrad/pastpapers/files/2019/paperia_1_2019.pdf}

\bibitem{tripos19p4}
University of Cambridge, Faculty of Mathematics.
\newblock \emph{Mathematical Tripos Part IA}, Paper~4, 5 de junio de 2019.
\newblock Preguntas 2E (series de términos $1/(n^{2}+n)$), 5E(b) (identidades de Fibonacci por inducción) y 7E(a) (caracterizaciones de la inyectividad), de Numbers and Sets.
\newblock \url{https://www.maths.cam.ac.uk/undergrad/pastpapers/files/2019/paperia_4_2019.pdf}

\bibitem{mit1801final}
Massachusetts Institute of Technology.
\newblock 18.01SC: \emph{Single Variable Calculus}, otoño de 2010. MIT OpenCourseWare.
\newblock Examen final con soluciones: problemas 3(a) (derivada a partir de la definición), 8(b) (integración por partes), 14 (serie de Taylor de $\ln(1+x)$ y cota del error) y~15 (desigualdades del arco tangente).
\newblock \url{https://ocw.mit.edu/courses/18-01sc-single-variable-calculus-fall-2010/pages/final-exam/}

\bibitem{mit1802ex2}
Massachusetts Institute of Technology.
\newblock 18.02SC: \emph{Multivariable Calculus}, otoño de 2010. MIT OpenCourseWare.
\newblock Unidad~2, «Partial Derivatives», examen~2 con soluciones: problemas~1 (gradiente, plano tangente, aproximación lineal y derivada direccional) y~3 (puntos críticos).
\newblock \url{https://ocw.mit.edu/courses/18-02sc-multivariable-calculus-fall-2010/pages/2.-partial-derivatives/exam-2/exam-materials/}

\bibitem{mit1805ex1}
Massachusetts Institute of Technology.
\newblock 18.05: \emph{Introduction to Probability and Statistics}, primavera de 2014. MIT OpenCourseWare.
\newblock Examen~1 con soluciones, problema~3 («Biased coins»).
\newblock \url{https://ocw.mit.edu/courses/18-05-introduction-to-probability-and-statistics-spring-2014/pages/exams/}

\bibitem{mit1805ex2}
Massachusetts Institute of Technology.
\newblock 18.05: \emph{Introduction to Probability and Statistics}, primavera de 2014. MIT OpenCourseWare.
\newblock Examen~2 con soluciones: problemas~2 (máxima verosimilitud para una Pareto) y~6 (contraste para la media de una normal).
\newblock \url{https://ocw.mit.edu/courses/18-05-introduction-to-probability-and-statistics-spring-2014/pages/exams/}

\subsection*{Material del curso y del proyecto}

\bibitem{apuntes}
Anatomía de un LLM.
\newblock \emph{Fase~0. Preliminares matemáticos}: apuntes de los módulos 0.1 a 0.9 (archivos P01 a P09), edición de octubre de 2026.
\newblock Cada módulo recoge en su bibliografía los libros de texto de referencia (Hammack, Lehman--Leighton--Meyer, Spivak, Apostol, Stewart, OpenStax, Blitzstein--Hwang, Rice y Wasserman, entre otros).

\bibitem{zipf}
G.~K. Zipf.
\newblock \emph{Human Behavior and the Principle of Least Effort}.
\newblock Addison-Wesley, 1949.

\bibitem{gutenberg}
Project Gutenberg.
\newblock Biblioteca de libros electrónicos de dominio público.
\newblock \url{https://www.gutenberg.org/}

\end{thebibliography}
